% ============================================================
% SCHEDE PG EXTRA — CASA GRIMALDI
% @rpg-schema: <https://rpg-schema.org/instances/genova1507/crew-grimaldi>
%   fitd:hasOptionalMember
%     <.../pg-bartolomeo-corso> ,
%     <.../pg-vanna-caruggi> .
% ============================================================

\chapter{Schede PG Extra --- Casa Grimaldi}
\index{schede PG!Grimaldi!extra}

% ════════════════════════════════════════════════════════════
% PG EXTRA 1 — BARTOLOMEO IL CORSO
% @rpg-schema: <https://rpg-schema.org/instances/genova1507/pg-bartolomeo-corso>
%   a fitd:Character ;
%   rdfs:label "Bartolomeo detto il Corso" ;
%   fitd:crew <.../crew-grimaldi> ;
%   fitd:role "Il Corsaro / Il Braccio Navale" ;
%   fitd:action [ fitd:actionName "Hunt" ; fitd:dots 3 ] ;
%   fitd:action [ fitd:actionName "Skirmish" ; fitd:dots 3 ] ;
%   fitd:action [ fitd:actionName "Wreck" ; fitd:dots 3 ] .
% ════════════════════════════════════════════════════════════

\pgheader[GrimaldiPurple]{Bartolomeo detto il Corso}{Il Corsaro / Il Braccio Navale}{Grimaldi}
\pgquote{Monaco paga bene e chiede poco. Finché dura, resto.}

\section*{Background \& Motivazione}

\begin{rulebox}{}
  \textbf{Background:} Trentanove anni, corso di nascita, corsaro per mestiere.
  Operate sotto patente dei Grimaldi da sei anni — il che lo rende tecnicamente
  un predatore autorizzato, non un pirata, distinzione a cui tiene molto.
  Ha una piccola flottiglia di due feluche basata a Monaco con cui si muove
  lungo la costa ligure e provenzale. È a Genova questa notte perché Rainier
  lo ha convocato con un messaggio breve: \textit{«Vieni. Potrebbe servire
  qualcuno che non ha paura dell'acqua.»}

  \medskip\noindent\textcolor{GrimaldiPurple}{\textbf{Motivazione:}} Bartolomeo
  lavora per denaro e per la libertà di movimento. I francesi hanno cominciato
  a fermare le sue feluche per ispezioni sistematiche. Questo non è accettabile.
  Chiunque liberi Genova dal presidio francese gli restituisce le rotte.
\end{rulebox}

\section*{Attributi \& Azioni}

\begin{multicols}{2}

\begin{attributoblock}[GrimaldiPurple]{INSIGHT --- Mente}
  \azione{Caccia}{Hunt}{3}{Rotte navali, posizionamento di navi, movimenti costieri}
  \azione{Studio}{Study}{1}{Navigazione, correnti, condizioni meteorologiche}
  \azione{Osserva}{Survey}{2}{Valuta situazioni in termini di rischio e fuga}
  \azione{Ingegno}{Tinker}{2}{Imbarcazioni, remi, vele, piccola artiglieria navale}
\end{attributoblock}


\begin{attributoblock}[GrimaldiPurple]{PROWESS --- Corpo}
  \azione{Finezza}{Finesse}{2}{Coltello, agilità su ponti scivolosi}
  \azione{Ombra}{Prowl}{2}{Sbarchi notturni, movimenti costieri non rilevati}
  \azione{Combattimento}{Skirmish}{3}{Violento e diretto — il mare non insegna eleganza}
  \azione{Forza bruta}{Wreck}{3}{Abbordaggi, sfondamenti, forza bruta marina}
\end{attributoblock}

\end{multicols}

\begin{attributoblock}[GrimaldiPurple]{RESOLVE --- Spirito}
  \begin{multicols}{2}
    \azione{Percezione}{Attune}{2}{Sente il pericolo in mare e sulla terraferma}
    \azione{Comando}{Command}{2}{Comanda il suo equipaggio con autorità pratica}
    \azione{Relazioni}{Consort}{2}{Corsari, capitani, taverne costiere}
    \azione{Persuasione}{Sway}{1}{Diretto. Se non funziona il discorso, cambia metodo.}
  \end{multicols}
\end{attributoblock}

\medskip
\begin{pgclockbox}{Stress \hfill \clock{9}{0}}
  \small\textit{Traumi: $\square$ Freddo \quad $\square$ Duro \quad $\square$ Ossessionato \quad $\square$ Paranoico \quad $\square$ Irrazionale \quad $\square$ Spezzato}
\end{pgclockbox}

\section*{Abilità Speciale}

% @rpg-schema: <.../pg-bartolomeo-corso>
%   fitd:specialAbility [ rdfs:label "Via di Mare" ] .

\begin{grimaldibox}{Via di Mare}
  Bartolomeo ha due feluche pronte a Sampierdarena. Può far uscire da Genova
  via mare fino a quattro persone in qualsiasi momento della sessione —
  senza tiri, senza tracce, senza domande. Il viaggio dura tre ore fino
  a Monaco.

  \medskip\noindent\textit{Sempre disponibile. Se viene usato durante
  un'azione di combattimento o inseguimento, richiede Tiro su Ombra.
  Le feluche non partono se il porto è sotto blocco navale attivo.}
\end{grimaldibox}

\section*{Legami}

\begin{table}[h]
\centering\renewcommand{\arraystretch}{1.4}
\begin{tabularx}{\linewidth}{@{}L{2.6cm} X@{}}
  \toprule\rowcolor{GrimaldiPurple}
  \textcolor{white}{\textbf{PG}} & \textcolor{white}{\textbf{Legame}} \\
  \midrule
  Rainier & Il signore mi paga, mi protegge e non mi chiede di rispettare
           regole che non capisco. È la definizione di buon datore di lavoro. \\
  \rowcolor{GrimaldiPurple!8}
  Ser Battista & È il tipo di soldato che capisce il mare senza averlo mai
                navigato davvero. Ci rispettiamo. \\
  Costantino & Sa dove si compra e si vende informazione su ogni porto
              del Mediterraneo. È il mio opposto e il mio complemento. \\
  \rowcolor{GrimaldiPurple!8}
  Serafina & Mi ha chiesto di insegnarle a stare su una barca. Rainier non
            lo sa. Sto valutando se è saggio darglielo a sapere. \\
  \bottomrule
\end{tabularx}
\end{table}

\section*{Equipaggiamento}
\begin{goldlist}
  \item Coltello da abbordaggio con impugnatura avvolta in corda marinara
  \item Mappa costiera della Liguria con i punti di sbarco non sorvegliati
  \item Patente di corsa firmata da Rainier II Grimaldi — lo distingue da un pirata
  \item Fischietto di segnalazione per l'equipaggio a Sampierdarena
\end{goldlist}

\clearpage

\begin{secretbox}
  \textbf{Segreto:} Una delle due feluche di Bartolomeo ha a bordo, in questo
  momento, un carico che non appartiene ai Grimaldi — merce di contrabbando
  che stava trasportando per conto di un mercante veneziano anonimo.
  Se la barca viene ispeziata — dai francesi o dai genovesi — il carico
  crea un problema diplomatico che potrebbe coinvolgere Monaco.
  Bartolomeo non ha ancora trovato il momento per dirlo a Rainier.

  \medskip\noindent\textcolor{SecretRed}{\textbf{Agenda:}}
  \textit{Far sparire il carico prima che qualcuno controlli le feluche.
  L'unico posto sicuro è il fondaco di Hamid nel sestiere di Soziglia —
  se riesce a portarcelo senza essere visto.}
\end{secretbox}

\clearpage

% ════════════════════════════════════════════════════════════
% PG EXTRA 2 — VANNA DEI CARUGGI
% @rpg-schema: <https://rpg-schema.org/instances/genova1507/pg-vanna-caruggi>
%   a fitd:Character ;
%   rdfs:label "Vanna dei Caruggi" ;
%   fitd:crew <.../crew-grimaldi> ;
%   fitd:role "La Contrabbandiera / La Radice Genovese" ;
%   fitd:action [ fitd:actionName "Survey" ; fitd:dots 3 ] ;
%   fitd:action [ fitd:actionName "Prowl" ; fitd:dots 4 ] ;
%   fitd:action [ fitd:actionName "Consort" ; fitd:dots 3 ] .
% ════════════════════════════════════════════════════════════

\pgheader[GrimaldiPurple]{Vanna dei Caruggi}{La Contrabbandiera / La Radice Genovese}{Grimaldi}
\pgquote{Monaco è lontana. I caruggi sono qui. Lavoro con quello che ho.}

\section*{Background \& Motivazione}

\begin{rulebox}{}
  \textbf{Background:} Trentatré anni, genovese pura, nata e cresciuta nel
  sestiere di Soziglia. Ha cominciato a fare la staffetta per i Grimaldi
  a dodici anni — portare messaggi tra il fondaco e il molo di Bartolomeo.
  Ora gestisce una rete di contrabbando che muove merci tra Monaco e i
  quartieri bassi di Genova senza che i doganieri francesi vedano niente.
  Non è un membro della famiglia nel senso formale — è la persona che conosce
  ogni vicolo, ogni porta secondaria, ogni guardia corruttibile tra Soziglia
  e il porto vecchio.

  \medskip\noindent\textcolor{GrimaldiPurple}{\textbf{Motivazione:}} Vanna
  odia i francesi per ragioni concrete: tasse sui mercati rionali, posti
  di blocco che rallentano il traffico, guardie che si comportano da
  occupanti. Non la interessa chi governerà Genova — le interessa che
  smettano le ispezioni. È disposta a rischiare molto per quella notte.
\end{rulebox}

\section*{Attributi \& Azioni}

\begin{multicols}{2}

\begin{attributoblock}[GrimaldiPurple]{INSIGHT --- Mente}
  \azione{Caccia}{Hunt}{2}{Trova persone, merci e passaggi in tutta la città}
  \azione{Studio}{Study}{0}{Ha imparato quello che serviva per strada}
  \azione{Osserva}{Survey}{3}{Legge le strade, le persone, le rotazioni di guardia}
  \azione{Ingegno}{Tinker}{2}{Nascondigli, camuffamenti di merci, vie non ufficiali}
\end{attributoblock}


\begin{attributoblock}[GrimaldiPurple]{PROWESS --- Corpo}
  \azione{Finezza}{Finesse}{2}{Borseggio, movimenti rapidi in spazi stretti}
  \azione{Ombra}{Prowl}{4}{I caruggi sono il suo habitat naturale}
  \azione{Combattimento}{Skirmish}{2}{Coltello, oggetti improvvisati, fuga rapida}
  \azione{Forza bruta}{Wreck}{1}{Non è il suo approccio preferito}
\end{attributoblock}

\end{multicols}

\begin{attributoblock}[GrimaldiPurple]{RESOLVE --- Spirito}
  \begin{multicols}{2}
    \azione{Percezione}{Attune}{2}{Sente quando un quartiere sta per esplodere}
    \azione{Comando}{Command}{1}{Guida i suoi contatti con pragmatismo, non autorità}
    \azione{Relazioni}{Consort}{3}{Bottegai, artigiani, gente di strada di tre sestieri}
    \azione{Persuasione}{Sway}{2}{Negozia veloce, parla la lingua giusta con le persone giuste}
  \end{multicols}
\end{attributoblock}

\medskip
\begin{pgclockbox}{Stress \hfill \clock{9}{0}}
  \small\textit{Traumi: $\square$ Freddo \quad $\square$ Duro \quad $\square$ Ossessionato \quad $\square$ Paranoico \quad $\square$ Irrazionale \quad $\square$ Spezzato}
\end{pgclockbox}

\section*{Abilità Speciale}

% @rpg-schema: <.../pg-vanna-caruggi>
%   fitd:specialAbility [ rdfs:label "La Città Sotto la Città" ] .

\begin{grimaldibox}{La Città Sotto la Città}
  Vanna conosce un percorso alternativo per qualsiasi destinazione
  all'interno delle mura di Genova — vicoli, cortili privati, passaggi
  sotterranei, tetti. Qualsiasi spostamento che lei guida \textbf{non può
  essere intercettato} da pattuglie o inseguitori, indipendentemente da
  quanti siano.

  \medskip\noindent\textit{Sempre attiva quando guida lei il gruppo.
  Se qualcuno del gruppo fa rumore o attira attenzione, il percorso
  viene compromesso e richiede Tiro su Ombra per continuare.}
\end{grimaldibox}

\section*{Legami}

\begin{table}[h]
\centering\renewcommand{\arraystretch}{1.4}
\begin{tabularx}{\linewidth}{@{}L{2.6cm} X@{}}
  \toprule\rowcolor{GrimaldiPurple}
  \textcolor{white}{\textbf{PG}} & \textcolor{white}{\textbf{Legame}} \\
  \midrule
  Rainier & Non l'ho mai incontrato di persona. Mi manda istruzioni tramite
           Costantino. Per me Monaco è un'astrazione — ma paga puntualmente. \\
  \rowcolor{GrimaldiPurple!8}
  Costantino & Il mio tramite. Lo rispetto perché non finge di essere
              qualcosa di diverso da quello che è. \\
  Bartolomeo & Il corso e io ci siamo divisi i ruoli senza discuterne:
              lui il mare, io la terra. Funziona. \\
  \rowcolor{GrimaldiPurple!8}
  Serafina & L'ho vista arrivare stamattina. Sembra la figlia di qualcuno
            che non ha mai avuto problemi veri. Probabilmente mi sbaglio.
            Probabilmente no. \\
  \bottomrule
\end{tabularx}
\end{table}

\section*{Equipaggiamento}
\begin{goldlist}
  \item Coltello corto con lama ceramica — non tintinna, non arrugginisce
  \item Mappa mentale completa dei caruggi di Soziglia, San Giovanni e Prè
  \item Lista di undici guardie francesi corruttibili con il prezzo di ognuna
  \item Chiave di un magazzino sotto il porto vecchio — accesso alla rete di contrabbando
\end{goldlist}

\clearpage

\begin{secretbox}
  \textbf{Segreto:} Vanna ha da tre giorni un ospite nascosto nel magazzino
  sotto il porto: un soldato francese disertore che ha rubato documenti
  dall'ufficio del governatore prima di scappare. Non sa cosa contengano
  i documenti — non sa leggere il francese. Li ha tenuti perché istintivamente
  ha capito che valgono qualcosa. Non lo ha detto ai Grimaldi perché voleva
  prima capire a chi conviene di più venderli.

  \medskip\noindent\textcolor{SecretRed}{\textbf{Agenda:}}
  \textit{Trovare qualcuno che legga i documenti e le dica cosa valgono
  davvero — e per chi. Costantino può farlo. Ma Costantino vende informazioni
  a tutti, e Vanna non sa ancora se questo è un problema o un'opportunità.}
\end{secretbox}
